\chapter{Host \& Network Security Auditing}
\label{ch:auditing}

Security auditing is not the same as penetration testing. Auditing checks whether
a system complies with security standards and policies. Penetration testing checks
whether it can be broken. Both are needed, and a pentester should understand where
auditing ends and testing begins.

% ─────────────────────────────────────────────────────────────────────────────
\section{What Security Auditing Is}
\label{sec:what_auditing}

\textbf{Definition:} Checking that everyone is following the rules and regulations
to secure systems. Auditing identifies risks to security --- anything that does not
comply with standards, regulations, or organisational policy.

\textbf{Why organisations do it:}
\begin{itemize}
  \item Risk management --- prioritise threats based on severity
  \item Identify loop holes before attackers do
  \item Build customer and regulatory confidence
  \item Compliance requirements (GDPR, HIPAA, ISO 27001, PCI-DSS)
  \item Ongoing process --- not a one-time event
\end{itemize}

% ─────────────────────────────────────────────────────────────────────────────
\section{Auditing vs Penetration Testing}
\label{sec:audit_vs_pentest}

\begin{notebox}
The relationship between the two: first conduct a security audit, then apply
the recommended remediations, then perform penetration testing to verify whether
the remediations actually fixed the problems.
\end{notebox}

\begin{longtable}{p{4cm} p{4cm} p{4cm}}
  \toprule
  \textbf{Aspect} & \textbf{Security Audit} & \textbf{Penetration Test} \\
  \midrule
  \textbf{Goal} & Compliance \& risk identification & Active exploitation \\
  \textbf{Method} & Review + interview + scan & Attack simulation \\
  \textbf{Output} & Compliance report & Vulnerability + exploit report \\
  \textbf{Timing} & Ongoing / periodic & After audit remediation \\
  \bottomrule
\end{longtable}

% ─────────────────────────────────────────────────────────────────────────────
\section{Audit Phases}
\label{sec:audit_phases}

\begin{enumerate}
  \item \textbf{Planning} --- define goals, understand existing documentation,
        schedule, scope
  \item \textbf{Information Gathering} --- number of computers, policies,
        employee access controls, encryption status, remote working risks
  \item \textbf{Field Work} --- actual audit execution, checking configurations,
        running vulnerability scans, interviewing staff
  \item \textbf{Analysis} --- compare findings against standards, identify
        improvement areas and weaknesses
  \item \textbf{Reporting} --- full findings report, improvement recommendations,
        gap analysis against the standard
\end{enumerate}

% ─────────────────────────────────────────────────────────────────────────────
\section{Types of Auditing}
\label{sec:audit_types}

\begin{itemize}
  \item \textbf{Access control audit} --- checking that only authorised persons
        access specific data and systems
  \item \textbf{Third-party audit} --- external auditors perform the assessment
  \item \textbf{Compliance audit} --- checking against government or industry
        standards (GDPR for EU data, HIPAA for patient data)
  \item \textbf{Technical audit} --- firewall rules, Linux hardening, network
        configuration, open ports
  \item \textbf{Network audit} --- open ports, security level, router testing
  \item \textbf{Web application audit} --- SQL injection, cross-site scripting,
        authentication flaws (using tools like Burp Suite)
\end{itemize}

% ─────────────────────────────────────────────────────────────────────────────
\section{NIST Framework}
\label{sec:nist}

NIST (National Institute of Standards and Technology) provides security
frameworks and guidelines that organisations use to develop security policies.
Understanding NIST codes is relevant because reports and remediation recommendations
often reference them.

\begin{notebox}
As a pentester, understanding NIST is important for writing reports that clients
and their compliance teams can act on. A finding mapped to a NIST control number
carries more weight than a plain description.
\end{notebox}

% ─────────────────────────────────────────────────────────────────────────────
\section{Lynis --- Linux System Auditing Tool}
\label{sec:lynis}

\tool{Lynis} performs local system audits --- it checks security hardening, installed
packages, kernel parameters, and configuration files. Used to audit the system
you are logged into.



\begin{termbox}
# Install
cd /opt/
gzip -d lynis-3.1.1.tar.gz
tar -xf lynis-3.1.1.tar
cd lynis/
chmod +x lynis

# Check system version
cat /etc/*issue

# Run a full system audit
./lynis audit system

# Run a specific test
./lynis audit system --tests HRDN-7230

# Run with auditor name in report
./lynis audit system --auditor "YourName"
\end{termbox}

\begin{notebox}

\begin{notebox}
Each Lynis finding includes a code (e.g., \ic{HRDN-7230}) and a remediation
suggestion. Click the code in the report for the full hardening recommendation.
\end{notebox}

% ─────────────────────────────────────────────────────────────────────────────
\section{Network Audit --- Practical Checklist}
\label{sec:network_audit}

\begin{termbox}
# Step 1 --- Network discovery
nmap -sn 192.168.x.0/24             # Live hosts

# Step 2 --- Port and service scan
nmap -sV -sC target_ip              # All services and versions

# Step 3 --- Vulnerability scanning
# Use Nessus (see Chapter 4) or MSF db_nmap

# Step 4 --- SSH authentication check
# Verify password-based auth is restricted
hydra -l root -P /usr/share/seclists/Passwords/xato-net-10-million-passwords.txt \
  ssh://target_ip:22 -t 2 -v
# If this finds a password, password auth is enabled and the password policy is weak
\end{termbox}

% ─────────────────────────────────────────────────────────────────────────────
\section{Penetration Test Report Structure}
\label{sec:pentest_report}

Every penetration test must produce a formal report. The structure I follow:

\begin{enumerate}
  \item \textbf{Executive Summary} --- high-level overview for non-technical readers
  \item \textbf{Methodology} --- tools used, phases followed
  \item \textbf{Findings} --- each vulnerability with severity, description,
        evidence (screenshots), and impact
  \item \textbf{Recommendations} --- specific steps to fix each finding
\end{enumerate}

\begin{takebox}
\textbf{Key Takeaways}
\begin{itemize}
  \item Auditing and penetration testing complement each other. Audit first,
        remediate, then pentest to verify.
  \item Lynis is the fastest way to get a security baseline on a Linux system.
  \item SSH with password authentication enabled and a weak password is a finding
        that will appear in almost every real audit. It comes up constantly.
  \item Understanding NIST matters for writing professional reports, not just
        running tools.
\end{itemize}
\end{takebox}

\end{notebox}